\documentclass[11pt,a4paper]{article}
% =====================================================================
%  CM3191 Organic Experiment 1 -- Report draft (pre-written)
%  Compile with XeLaTeX. Set \draftfalse to hide the Korean notes page
%  and the grey placeholder boxes' labels before submission.
% =====================================================================
\newif\ifdraft \drafttrue

\usepackage[a4paper,margin=22mm]{geometry}
\usepackage{fontspec}
\usepackage{xeCJK}
\setmainfont{TeX Gyre Termes}
\setsansfont{TeX Gyre Heros}
\setCJKmainfont{Noto Serif CJK KR}
\xeCJKsetup{CJKspace=true, CJKecglue={}}
\usepackage{amsmath,amssymb}
\usepackage[version=4]{mhchem}
\usepackage{chemfig}
\usepackage{siunitx}
\sisetup{detect-all, per-mode=symbol}
\usepackage{booktabs,tabularx,array}
\usepackage{enumitem}
\usepackage{xcolor}
\usepackage[super]{cite}
\usepackage{titlesec}
\usepackage[hidelinks]{hyperref}
\usepackage{caption}
\usepackage{float}
\captionsetup{font=small,labelfont=bf}
\titleformat{\section}{\large\bfseries}{\thesection.}{0.5em}{}
\titleformat{\subsection}{\normalsize\bfseries}{\thesubsection}{0.5em}{}
\setlength{\parskip}{3pt}
\linespread{1.1}
\setchemfig{atom sep=1.9em, bond style={line width=0.6pt}}

% ---- placeholders: everything in a grey box must be replaced after the lab
\newcommand{\res}[1]{\fcolorbox{black}{gray!18}{\footnotesize\sffamily #1}}
\newcommand{\obs}[1]{\fcolorbox{black}{gray!18}{\footnotesize\sffamily OBS: #1}}
\newcommand{\resp}[1]{\par\smallskip\noindent\fcolorbox{black}{gray!18}{\parbox{\dimexpr\linewidth-2\fboxsep-2\fboxrule\relax}{\footnotesize\sffamily #1}}\par\smallskip}

\begin{document}

% =====================================================================
\ifdraft
\thispagestyle{empty}
\begin{center}{\Large\bfseries 작성 안내 (제출 전 이 페이지 삭제)}\end{center}
\vspace{2mm}
\textbf{이 초안의 구조.} 결과와 무관한 부분(서론·이론, 실험 방법, 고찰의 이론적 틀, 참고문헌)은 완성된 문장으로 써 두었습니다. 실험 후 바뀌는 값과 관찰은 \res{회색 상자}로 남겼습니다. 상자 안의 숫자는 매뉴얼의 계획값이라 컴파일해도 문서가 자연스럽게 읽히지만, \textbf{반드시 본인 측정값으로 바꿔야 합니다}.

\textbf{실험 후 채울 것}
\begin{enumerate}[leftmargin=6mm,itemsep=0.5mm]
\item 2.2절과 표 1: 실제로 잰 질량과 mmol, 당량(datasheet 사진 기준). 재결정에 쓴 에탄올 양.
\item 3절: 생성물 질량, 수율, TON/TOF(표 2), 관찰(표 3), IR·NMR 실측값(표 4, 5).
\item 4절: \texttt{\%} 주석으로 달아 둔 ``결과가 A이면 / B이면'' 문장 중 맞는 쪽 선택.
\item 5절 결론의 숫자, 부록(서명된 datasheet, 스펙트럼 전체와 확대).
\end{enumerate}

\textbf{제출 전 확인할 것 [CONFIRM]}
\begin{enumerate}[leftmargin=6mm,itemsep=0.5mm]
\item 유기 파트 보고서 형식·분량·마감. 일정표에 ``평가 방식은 7주차 전 업데이트''로만 되어 있습니다. proforma 형식이면 이 초안의 문단을 답안으로 옮겨 쓰면 됩니다.
\item 과목의 생성형 AI 사용 규정. 규정에 따라 사용 사실 표기가 필요할 수 있습니다.
\item NMR 용매(DMSO-$d_6$로 가정)와 측정 주파수(400 MHz로 가정), IR 기기명.
\item 조교 서명 datasheet와 스펙트럼을 부록으로 붙이는지.
\end{enumerate}
\textbf{컴파일}: XeLaTeX. 첫 줄 근처의 \verb|\drafttrue|를 \verb|\draftfalse|로 바꾸면 이 페이지가 빠집니다.
\newpage
\setcounter{page}{1}
\fi
% =====================================================================

\begin{center}
{\LARGE\bfseries Aqueous Suzuki--Miyaura Cross-Coupling:\\[1mm] Synthesis of 4-Phenylbenzoic Acid}\\[4mm]
{\large Kim Donghwan (A0266094X)}\\[1mm]
CM3191 Chemical Synthesis Experiments, Organic Experiment 1\\
Thursday laboratory, fume hood B7 \quad\textbullet\quad Date of experiment: 1 October 2026\\
Instructor: Dr Hoang Truong Giang
\end{center}
\vspace{2mm}
\noindent\rule{\linewidth}{0.6pt}

\noindent\textbf{Abstract.} 4-Phenylbenzoic acid was prepared by the palladium-catalysed Suzuki--Miyaura coupling of 4-bromobenzoic acid with phenylboronic acid in water, using aqueous \ce{Na2CO3} as the base and 0.015 mol\% of a water-soluble Pd complex of 2-amino-4,6-dihydroxypyrimidine at 70\,\si{\celsius} for 30 min. The product was isolated by acidification, recrystallised from ethanol and obtained as a \res{white} solid in \res{XX}\% yield. Its identity was confirmed by \textsuperscript{1}H NMR and IR spectroscopy. The isolated yield corresponds to a turnover number of \res{XXXX} and an average turnover frequency of \res{XXX}~\si{\per\minute}.

\noindent\rule{\linewidth}{0.6pt}

% =====================================================================
\section{Introduction}

Palladium-catalysed cross-coupling reactions join two organic fragments by assembling them on a metal centre, where they form a new carbon--carbon single bond. The Suzuki--Miyaura reaction couples an organoboron compound with an organic halide or triflate.\cite{Miyaura1979,Miyaura1995} It was recognised, together with the Heck and Negishi couplings, by the 2010 Nobel Prize in Chemistry.\cite{Nobel2010} In a survey of reactions used in medicinal chemistry, the Suzuki--Miyaura coupling accounted for about 40\% of all C--C bond-forming reactions.\cite{Roughley2011} Its popularity comes from the properties of the boron reagents: boronic acids are generally bench-stable solids, tolerate many functional groups and have low toxicity, and their boron by-products are easily removed in aqueous waste.\cite{Lennox2014}

The generally accepted mechanism is a Pd(0)/Pd(II) cycle (Figure~\ref{sch:cycle}). \emph{Oxidative addition} of the aryl halide to Pd(0) gives an arylpalladium(II) halide. The rate of this step follows the order Ar--I $>$ Ar--Br $\approx$ Ar--OTf $\gg$ Ar--Cl, reflecting the increasing strength of the carbon--leaving group bond. In \emph{transmetalation}, the aryl group on boron is transferred to palladium. This step requires a base, because the neutral, largely covalent boronic acid does not transmetalate readily. Two pathways have been proposed: the base either converts the boronic acid into a more nucleophilic trihydroxyborate (the ``boronate'' pathway), or it replaces the halide on palladium to give an arylpalladium hydroxide that reacts with the neutral boronic acid (the ``oxo-palladium'' pathway). Stoichiometric rate measurements by Carrow and Hartwig indicate that the oxo-palladium pathway dominates under weak-base, aqueous conditions,\cite{Carrow2011} although the preferred route may depend on the conditions. Finally, \emph{reductive elimination} forms the biaryl C--C bond and regenerates Pd(0).

\begin{figure}[H]
\centering
\setchemfig{atom sep=1.6em}
\schemestart
\chemfig{HOOC-*6(-=-(-Br)=-=)}
\+
\chemfig{*6(-=-(-B(-[:30]OH)(-[:-30]OH))=-=)}
\arrow{->[\footnotesize catalyst \textbf{1} (0.015 mol\%)][\footnotesize \ce{Na2CO3}, \ce{H2O}, 70\,\si{\celsius}, 30 min]}[,2.3]
\chemfig{HOOC-*6(-=-(-*6(-=-=-=))=-=)}
\schemestop
\par\vspace{3mm}
{\small $\ce{Pd^0L_n} \xrightarrow{\text{Ar--Br}} \ce{Ar-Pd^{II}L_n-Br} \xrightarrow{\ce{OH^-}} \ce{Ar-Pd^{II}L_n-OH} \xrightarrow{\ce{PhB(OH)2}} \ce{Ar-Pd^{II}L_n-Ph} \xrightarrow{-\,\text{Ar--Ph}} \ce{Pd^0L_n}$}
\captionof{figure}{Reaction studied in this experiment (top) and the simplified catalytic cycle, drawn here via the oxo-palladium pathway (bottom). Ar = 4-carboxylatophenyl; catalyst \textbf{1} = Pd(OAc)\textsubscript{2}/2-amino-4,6-dihydroxypyrimidine (1:2) in aqueous NaOH.}
\label{sch:cycle}
\end{figure}

Traditional Suzuki couplings are often run in organic solvents such as toluene or DMF. Replacing them with water follows the green chemistry principle of using safer solvents,\cite{Anastas1998} but it raises two problems: most organic substrates and conventional Pd--phosphine catalysts do not dissolve in water. The procedure used here, adapted from Hamilton \emph{et al.},\cite{Hamilton2013} solves both. First, under the basic conditions both coupling partners are anions (a carboxylate and a boronate) and are therefore water-soluble. Second, the catalyst is a phosphine-free Pd complex of 2-amino-4,6-dihydroxypyrimidine that is soluble in water and active at very low loading.\cite{Chalker2009} In addition, the product is a carboxylic acid, so it can be isolated simply by acidifying the reaction mixture, without any extraction with organic solvents.

The aims of this experiment were (i) to synthesise 4-phenylbenzoic acid by an aqueous Suzuki--Miyaura coupling, (ii) to purify it by recrystallisation from ethanol, (iii) to confirm its identity and purity by IR and \textsuperscript{1}H NMR spectroscopy, and (iv) to evaluate the yield and the catalyst performance in terms of turnover number (TON) and turnover frequency (TOF).

% =====================================================================
\section{Experimental}

\subsection{General}
4-Bromobenzoic acid, phenylboronic acid, aqueous \ce{Na2CO3} (0.5 M), hydrochloric acid (1 M) and ethanol were used as supplied. The palladium catalyst solution (0.25 mM in Pd) was prepared by the laboratory staff: \ce{Pd(OAc)2} (5.6 mg, 0.025 mmol) and 2-amino-4,6-dihydroxypyrimidine (6.3 mg, 0.050 mmol) were treated with a solution of NaOH (4.0 mg, 0.10 mmol) in deionised water (10 mL) at 60\,\si{\celsius} for 20 min, and the resulting solution was diluted to 100 mL.\cite{Manual2026} IR spectra were recorded on \res{instrument} with an ATR accessory. \textsuperscript{1}H NMR spectra were recorded at \res{400} MHz in \res{DMSO-$d_6$}; chemical shifts ($\delta$) are reported in ppm relative to the residual solvent signal (DMSO-$d_5$, $\delta$ 2.50).\cite{Fulmer2010}

\subsection{Synthesis of 4-phenylbenzoic acid}
4-Bromobenzoic acid (\res{0.340} g, \res{1.69} mmol, 1.00 equiv) and phenylboronic acid (\res{0.250} g, \res{2.05} mmol, \res{1.21} equiv) were placed in a 20 mL reaction vial containing a magnetic stirrer bar. Aqueous \ce{Na2CO3} (0.5 M, 10 mL, 5.0 mmol, \res{2.96} equiv) was added, and the capped vial was stirred vigorously at room temperature until all of the solid had dissolved. The palladium catalyst solution (0.25 mM, 1.0 mL, 0.25 µmol Pd, \res{0.015} mol\%) was then added, and the capped vial was heated in a preheated aluminium block at 70\,\si{\celsius} for 30 min. \obs{appearance before/after heating}

The mixture was cooled to room temperature, transferred to a 250 mL conical flask and stirred in an ice bath while hydrochloric acid (1 M, 20 mL) was added dropwise. Gas evolution was observed \obs{extent} and a \res{white} precipitate formed. After the mixture had been confirmed to be acidic with pH paper, water (20 mL) was added and stirring was continued in the ice bath for 10 min. The solid was collected by suction filtration.

\subsection{Recrystallisation}
The crude solid was transferred to a 50 mL conical flask with a stirrer bar, and 1 M HCl (3 mL) and ethanol (10 mL) were added. The mixture and a separate flask of ethanol (20 mL, with boiling chips) were heated to a gentle boil, and the hot ethanol was added dropwise until all of the solid had dissolved (\res{X} mL in total). The solution was allowed to cool undisturbed to room temperature (10--15 min) and was then cooled in an ice bath for 15 min, together with the remaining ethanol. The crystals were collected by suction filtration, washed with a minimal volume of cold ethanol and dried under an IR lamp for 15 min to give 4-phenylbenzoic acid as a \res{white crystalline} solid (\res{0.XXX} g, \res{X.XX} mmol, \res{XX}\%).

\noindent\textbf{Characterisation data.} IR (ATR) $\tilde{\nu}_\text{max}$/\si{\per\centi\metre}: \res{values}. \textsuperscript{1}H NMR (\res{400} MHz, \res{DMSO-$d_6$}) $\delta$: \res{values}.

% =====================================================================
\section{Results}

\begin{table}[H]
\centering
\caption{Quantities of reagents (planned values from the laboratory manual\cite{Manual2026} and actual values used).}
\label{tab:reagents}
\small
\begin{tabular}{lccccc}
\toprule
Reagent & $M$ / conc. & Planned & Actual & $n$ / mmol & Equiv. \\
\midrule
4-Bromobenzoic acid & 201.02 g/mol & 0.34 g & \res{} & \res{1.69} & 1.00 \\
Phenylboronic acid & 121.93 g/mol & 0.25 g & \res{} & \res{2.05} & \res{1.21} \\
\ce{Na2CO3} (aq) & 0.5 M & 10 mL & \res{} & 5.0 & \res{2.96} \\
Pd catalyst solution & 0.25 mM & 1.0 mL & \res{} & 0.00025 & \res{0.00015} \\
\midrule
4-Phenylbenzoic acid & 198.22 g/mol & theor.\ \res{0.335} g & \res{} & \res{} & --- \\
\bottomrule
\end{tabular}
\end{table}

\begin{table}[H]
\centering
\caption{Yield and catalyst performance.}
\label{tab:yield}
\small
\begin{tabular}{lc}
\toprule
Quantity & Value \\
\midrule
Mass of purified product & \res{0.XXX} g \\
Theoretical yield (from actual mass of 4-bromobenzoic acid) & \res{0.XXX} g \\
Isolated yield & \res{XX}\% \\
TON $= n_\text{product}/n_\text{Pd}$ \ ($n_\text{Pd} = 2.5\times10^{-4}$ mmol) & \res{XXXX} \\
TOF $=$ TON$/t$ \ ($t = 30$ min) & \res{XXX} \si{\per\minute} \ (\res{XXXX} \si{\per\hour}) \\
\bottomrule
\end{tabular}
\end{table}

\begin{table}[H]
\centering
\caption{Observations during the experiment.}
\label{tab:obs}
\small
\begin{tabularx}{\linewidth}{lX}
\toprule
Stage & Observation \\
\midrule
Solids in \ce{Na2CO3} solution & \obs{dissolution, colour, time} \\
After addition of catalyst and 30 min at 70\,\si{\celsius} & \obs{colour, any solid or grey tint} \\
Acidification with 1 M HCl & \obs{gas, precipitate, pH paper colour} \\
Recrystallisation & \obs{ethanol volume, crystal formation, appearance} \\
Dried product & \obs{colour, form} \\
\bottomrule
\end{tabularx}
\end{table}

\begin{table}[H]
\centering
\caption{\textsuperscript{1}H NMR data of the product compared with literature values (DMSO-$d_6$). Labels a--e refer to the structure in the Discussion.}
\label{tab:nmr}
\small
\begin{tabular}{lccccl}
\toprule
H & Lit.\ $\delta$\cite{CBNMR} & Obs.\ $\delta$ & Mult. & Integral & Assignment \\
\midrule
COOH & 13.05 & \res{} & br s & \res{} & carboxylic acid \\
a & 8.07 & \res{} & d & \res{} & central ring, \emph{ortho} to \ce{CO2H} \\
b & 7.82 & \res{} & d & \res{} & central ring, \emph{ortho} to Ph \\
c & 7.75 & \res{} & d & \res{} & terminal ring, \emph{ortho} \\
d & 7.52 & \res{} & t & \res{} & terminal ring, \emph{meta} \\
e & 7.44 & \res{} & t & \res{} & terminal ring, \emph{para} \\
\bottomrule
\end{tabular}
\end{table}

\begin{table}[H]
\centering
\caption{Principal IR absorptions of the product.}
\label{tab:ir}
\small
\begin{tabular}{lcc}
\toprule
Vibration & Typical range\cite{Pavia} / \si{\per\centi\metre} & Observed / \si{\per\centi\metre} \\
\midrule
O--H stretch (H-bonded carboxylic acid dimer) & 3300--2500 (broad) & \res{} \\
C=O stretch (aryl carboxylic acid) & ca.\ 1710--1680 & \res{} \\
Aromatic C=C stretch & ca.\ 1600, 1500--1450 & \res{} \\
C--O stretch & 1320--1210 & \res{} \\
C--H out-of-plane bend, \emph{para}-disubstituted ring & ca.\ 860--800 & \res{} \\
C--H out-of-plane bend, monosubstituted ring & ca.\ 770--730 and 710--690 & \res{} \\
\bottomrule
\end{tabular}
\end{table}

% =====================================================================
\section{Discussion}

\subsection{Why the coupling works in water}
Three features of the system make water a suitable solvent. First, the substrate carries a carboxylic acid group, which acts as a solubilising ``handle''. Aqueous \ce{Na2CO3} (0.5 M, estimated pH $\approx$ 12 from $K_\text{b}(\ce{CO3^{2-}}) = K_\text{w}/K_\text{a2} = 2.1\times10^{-4}$) deprotonates it completely: with a p$K_\text{a}$ close to that of benzoic acid (4.2), the carboxylate-to-acid ratio is about $10^{12-4.2}\approx6\times10^{7}$. At the same pH the boronic acid is largely converted into its borate. Both partners are therefore in homogeneous solution, which is why the catalyst was added only after all solids had dissolved. Second, the pyrimidine ligand, deprotonated by NaOH during catalyst preparation, makes the Pd complex water-soluble without phosphines, so the reaction does not need to be protected from air.\cite{Chalker2009} Third, the base plays several roles at once: it solubilises the substrate, enables transmetalation (Section~1) and neutralises the acidic by-products. The overall process under the reaction conditions can be written as
\[
\ce{Br-C6H4-CO2^- + PhB(OH)2 + OH^- -> Ph-C6H4-CO2^- + B(OH)3 + Br^-}
\]
which is balanced in atoms and charge.

The main limitations of the aqueous method are that it relies on a substrate with an ionisable group and on a specially designed water-soluble catalyst. In addition, boronic acids can undergo protodeboronation in hot basic water, which is one reason for using a slight excess (1.2 equiv) of phenylboronic acid.

\subsection{Acid--base control of isolation and purification}
The isolation depends entirely on the carboxylic acid being switched between its soluble salt and its insoluble neutral form. Hydrochloric acid (20 mmol) first destroys the carbonate ($\ce{CO3^{2-} + 2H+ -> CO2 + H2O}$, which accounts for the gas evolution) and then protonates the product. Because 5.0 mmol of carbonate consumes 10 mmol of acid, about 10 mmol of HCl remains in excess in roughly 51 mL of solution. This gives $[\ce{H+}] \approx 0.20$ M and pH $\approx$ 0.7, at which only about $10^{0.7-4.2}\approx3\times10^{-4}$ of the product remains as the carboxylate. Cooling in ice lowers its solubility further and absorbs the heat of neutralisation. Dilution with water then keeps the inorganic salts and boron by-products in solution, so they pass into the filtrate.

In the recrystallisation, a small amount of 1 M HCl was added to keep the product fully in its neutral acid form. Using the minimum volume of hot ethanol is important because some product always remains dissolved in the cold mother liquor. By-products such as biphenyl (from homocoupling of phenylboronic acid) are expected to stay in the ethanol mother liquor. Unreacted 4-bromobenzoic acid, however, is chemically similar to the product and may co-crystallise, so purity must be checked by NMR. \resp{Discuss your ethanol volume (X mL) and whether crystallisation was slow or rapid.}

\subsection{Identity and purity of the product}
\begin{center}
\setchemfig{atom sep=2.1em}
\chemfig{HOOC-*6(-(-[,0.45,,,draw=none]{\small a})=(-[,0.45,,,draw=none]{\small b})-(-*6(-(-[,0.45,,,draw=none]{\small c})=(-[,0.45,,,draw=none]{\small d})-(-[,0.45,,,draw=none]{\small e})=-=))=-=)}
\end{center}
The most informative evidence for the coupling is the appearance of the monosubstituted phenyl ring. 4-Bromobenzoic acid shows only two two-proton doublets in the aromatic region, whereas the product also shows the terminal-ring protons c--e, including two triplets (d, 2H; e, 1H) that cannot arise from the starting material. The \textsuperscript{1}H NMR spectrum of the product showed \res{signals at \dots}, in \res{good} agreement with the literature values in Table~\ref{tab:nmr}. The total aromatic integral of \res{X.XX} relative to H\textsubscript{e} = 1.00 \res{matches / exceeds} the expected 9H.
% --- choose one of the following sentences after the lab ---
% (A) No signals attributable to unreacted 4-bromobenzoic acid were observed, indicating that the product was obtained in high purity.
% (B) Additional doublets at delta X.XX and X.XX indicate residual 4-bromobenzoic acid, estimated at X mol% from the integrals relative to H_e.
% (C) Signals at delta 1.06 and 3.44 show residual ethanol, i.e. the product was not completely dry; this would inflate the isolated yield.

The IR spectrum (Table~\ref{tab:ir}) showed the broad O--H absorption and the C=O stretch typical of a hydrogen-bonded aryl carboxylic acid. These bands are also present in the starting material, so they confirm that the carboxylic acid group survived but do not by themselves prove the coupling. The out-of-plane C--H bending bands near \res{750} and \res{690}~\si{\per\centi\metre}, characteristic of a monosubstituted benzene ring, provide supporting evidence for the new phenyl group.

\subsection{Yield}
The isolated yield was \res{XX}\%. Because every operation after the reaction depends on precipitation and recrystallisation, losses to the filtrates are expected, especially to the ethanol mother liquor.
% --- keep the points that apply ---
% * Low yield: excess hot ethanol (X mL used); transfer losses on filter paper and glassware; incomplete precipitation if the mixture was not acidic enough or not cold enough; incomplete reaction (check NMR for starting material).
% * Yield near or above 100%: incomplete drying (ethanol/water in NMR) or inorganic salts (NaCl) carried into the product.
\resp{Explain your yield using the datasheet observations.}

\subsection{Catalyst performance and green chemistry}
With only $2.5\times10^{-4}$ mmol Pd (0.015 mol\%, 148 ppm relative to the aryl bromide), each palladium centre turned over on average \res{XXXX} times in 30 min (TON), corresponding to a mean TOF of \res{XXX}~\si{\per\minute}. These values are lower limits for three reasons: they are based on the isolated rather than the conversion yield, they assume that all of the added palladium was catalytically active, and they average over the full 30 min even if the reaction finished earlier.

From a green chemistry perspective, the reaction uses water as the only reaction solvent, requires no inert atmosphere and needs no extraction solvent. Its atom economy, calculated from the two coupling partners only,\cite{Trost1991} is $198.22/(201.02+121.93) = 61.4\%$; the remainder leaves mainly as boric acid and bromide salts. The main waste streams are the acidic aqueous filtrate, which contains the salts and traces of palladium and must be collected as metal-containing waste, and the ethanol mother liquor. Because a large mass of water is used, the mass-based E-factor would be dominated by water; it is common practice to report it both with and without water.\cite{Sheldon2007}

% =====================================================================
\section{Conclusion}
4-Phenylbenzoic acid was synthesised by an aqueous Suzuki--Miyaura coupling of 4-bromobenzoic acid and phenylboronic acid, using a water-soluble Pd catalyst at 0.015 mol\%. It was isolated by acidification and purified by recrystallisation from ethanol in \res{XX}\% yield. \textsuperscript{1}H NMR and IR spectroscopy confirmed its identity, in particular through the signals of the new monosubstituted phenyl ring, and \res{no significant / X mol\%} starting material remained. The catalyst achieved a TON of \res{XXXX} and a TOF of \res{XXX}~\si{\per\minute}. The experiment shows how a suitable choice of substrate, base and catalyst allows a cross-coupling reaction to be carried out and worked up in water.

% =====================================================================
\begin{thebibliography}{99}
\small
\setlength{\itemsep}{1pt}
\bibitem{Miyaura1979} Miyaura, N.; Yamada, K.; Suzuki, A. A New Stereospecific Cross-Coupling by the Palladium-Catalyzed Reaction of 1-Alkenylboranes with 1-Alkenyl or 1-Alkynyl Halides. \emph{Tetrahedron Lett.} \textbf{1979}, \emph{20}, 3437--3440.
\bibitem{Miyaura1995} Miyaura, N.; Suzuki, A. Palladium-Catalyzed Cross-Coupling Reactions of Organoboron Compounds. \emph{Chem. Rev.} \textbf{1995}, \emph{95}, 2457--2483.
\bibitem{Nobel2010} The Nobel Prize in Chemistry 2010. NobelPrize.org. \url{https://www.nobelprize.org/prizes/chemistry/2010/summary/} (accessed \res{date}).
\bibitem{Roughley2011} Roughley, S.\,D.; Jordan, A.\,M. The Medicinal Chemist's Toolbox: An Analysis of Reactions Used in the Pursuit of Drug Candidates. \emph{J. Med. Chem.} \textbf{2011}, \emph{54}, 3451--3479.
\bibitem{Lennox2014} Lennox, A.\,J.\,J.; Lloyd-Jones, G.\,C. Selection of Boron Reagents for Suzuki--Miyaura Coupling. \emph{Chem. Soc. Rev.} \textbf{2014}, \emph{43}, 412--443.
\bibitem{Carrow2011} Carrow, B.\,P.; Hartwig, J.\,F. Distinguishing Between Pathways for Transmetalation in Suzuki--Miyaura Reactions. \emph{J. Am. Chem. Soc.} \textbf{2011}, \emph{133}, 2116--2119.
\bibitem{Anastas1998} Anastas, P.\,T.; Warner, J.\,C. \emph{Green Chemistry: Theory and Practice}; Oxford University Press: Oxford, 1998.
\bibitem{Hamilton2013} Hamilton, A.\,E.; Buxton, A.\,M.; Peeples, C.\,J.; Chalker, J.\,M. An Operationally Simple Aqueous Suzuki--Miyaura Cross-Coupling Reaction for an Undergraduate Organic Chemistry Laboratory. \emph{J. Chem. Educ.} \textbf{2013}, \emph{90}, 1509--1513.
\bibitem{Chalker2009} Chalker, J.\,M.; Wood, C.\,S.\,C.; Davis, B.\,G. A Convenient Catalyst for Aqueous and Protein Suzuki--Miyaura Cross-Coupling. \emph{J. Am. Chem. Soc.} \textbf{2009}, \emph{131}, 16346--16347.
\bibitem{Manual2026} Hoang, T.\,G. \emph{CM3191 Chemical Synthesis Experiments: Organic Lab Manual}; National University of Singapore, AY2026/27; Experiment 1, pp 2--7.
\bibitem{Fulmer2010} Fulmer, G.\,R.; Miller, A.\,J.\,M.; Sherden, N.\,H.; Gottlieb, H.\,E.; Nudelman, A.; Stoltz, B.\,M.; Bercaw, J.\,E.; Goldberg, K.\,I. NMR Chemical Shifts of Trace Impurities: Common Laboratory Solvents, Organics, and Gases in Deuterated Solvents Relevant to the Organometallic Chemist. \emph{Organometallics} \textbf{2010}, \emph{29}, 2176--2179.
\bibitem{CBNMR} 4-Biphenylcarboxylic acid (CAS 92-92-2), \textsuperscript{1}H NMR spectrum (400 MHz). ChemicalBook. \url{https://www.chemicalbook.com/SpectrumEN_92-92-2_1HNMR.htm} (accessed \res{date}).
\bibitem{Pavia} Pavia, D.\,L.; Lampman, G.\,M.; Kriz, G.\,S.; Vyvyan, J.\,R. \emph{Introduction to Spectroscopy}, 5th ed.; Cengage Learning: Stamford, CT, 2015.
\bibitem{Trost1991} Trost, B.\,M. The Atom Economy---A Search for Synthetic Efficiency. \emph{Science} \textbf{1991}, \emph{254}, 1471--1477.
\bibitem{Sheldon2007} Sheldon, R.\,A. The E Factor: Fifteen Years On. \emph{Green Chem.} \textbf{2007}, \emph{9}, 1273--1283.
\end{thebibliography}

\end{document}
